% titel: Blatt 8, A6: Dipolmoment OF2/BeF2 und PF3/BF3
% skript: Mortimer Kap. 4 (Molekülgeometrie, Dipolmoment)

% Teil I: Werkzeuge, kein Lösungsweg.

Ein Molekül hat nur dann ein permanentes Dipolmoment $\mu\neq 0$, wenn die
Bindungsdipole sich \emph{nicht} zu Null addieren. Die Geometrie entscheidet
darüber.

\begin{werkzeugkasten}[Dipolmoment und Symmetrie]
Das Dipolmoment ist eine \emph{Vektorgrösse}:
\[
  \vec\mu = \sum_i \vec\mu_i\ (\text{alle Bindungs- und freien
  Elektronenpaarbeiträge}).
\]
\begin{itemize}
  \item \textbf{Symmetrische} Moleküle (linear, trigonal-planar, mit gleichen
        Liganden) heben sich die Teilmomente gegenseitig auf $\Rightarrow
        \mu = 0$.
  \item \textbf{Nicht-symmetrische} Geometrien (gewinkelt, pyramidal oder
        mit verschiedenen Atomen) ergeben ein Restmoment $\Rightarrow
        \mu\neq 0$.
\end{itemize}
\textit{Skript:} Mortimer Kap. 4 (Form: VSEPR; Dipolmoment).
\end{werkzeugkasten}

\begin{achtung}[Darauf achten]
\begin{itemize}
  \item Frage ist: \emph{Molekül}geometrie, nicht nur Bindungspolarität.
        BeF$_2$ und BF$_3$ haben polare Bindungen, aber die Geometrie macht
        das Gesamtmoment Null.
  \item OF$_2$ ist (wie H$_2$O) gewinkelt; PF$_3$ ist pyramidal. Beide haben
        deshalb ein Restmoment $\neq 0$.
  \item Nicht das Vorhandensein eines freien Paars allein, sondern dessen
        Umgebung (Symmetrie) bestimmt das Dipolmoment.
\end{itemize}
\end{achtung}
